% Generated draft: requires mathematical and source review.
\documentclass[11pt]{article}
\usepackage[margin=25mm]{geometry}
\usepackage{fontspec}
\setmainfont{Latin Modern Roman}
\usepackage{amsmath,amssymb,mathtools,unicode-math}
\setmathfont{Latin Modern Math}
\usepackage{enumitem,needspace,xurl,hyperref}
\setlength{\parindent}{0pt}
\setlength{\parskip}{5pt}
\setlength{\emergencystretch}{3em}
\newcommand{\N}{\mathbb{N}}
\newcommand{\Z}{\mathbb{Z}}
\newcommand{\Q}{\mathbb{Q}}
\newcommand{\R}{\mathbb{R}}
\newcommand{\C}{\mathbb{C}}
\newcommand{\F}{\mathbb{F}}
\newcommand{\A}{\mathbb{A}}
\newcommand{\PP}{\mathbb{P}}
\newcommand{\E}{\mathbb{E}}
\newcommand{\eps}{\varepsilon}
\newcommand{\dd}{\,\mathrm{d}}
\begin{document}
\section*{UnsolvedMath ID: 30000678}
\textbf{Source problem code:} catalogue entry
\subsection{Definitions and mathematical statement}
Let \(R\) be an associative unital ring, not assumed commutative, and work with unital left modules. A left module \(C\) is cotorsion if
\[
\operatorname{Ext}_R^1(F,C)=0\qquad\text{for every flat left }R\text{-module }F.
\]
An exact sequence of left modules is pure if it remains exact after tensoring with every right module. A left module \(E\) is pure-injective if maps into \(E\) extend across pure monomorphisms. The first source question asks for a ring \(R\) that is noncoherent but satisfies
\[
\forall C\in R\text{-}\operatorname{Mod},\qquad
C\text{ cotorsion}\ \Longrightarrow\ C\text{ pure-injective}.
\]
The quantifier is over all left modules, not only finitely generated ones. For a noncommutative ring, right coherence means every finitely generated right ideal is finitely presented, and left coherence is the analogous left-sided condition. The source does not specify which coherence fails; a side cannot be selected silently from the left-module wording.

For a left module \(M\), let \(i_M:M\hookrightarrow\operatorname{PE}(M)\) be its pure-injective envelope, with its usual pure embedding and minimal-envelope property, and put
\[
C_M=\operatorname{coker}i_M=\operatorname{PE}(M)/i_M(M).
\]
The second question asks for ring-theoretic properties equivalent to
\[
\forall M\in R\text{-}\operatorname{Mod},\qquad C_M\text{ is pseudoflat}.
\]
The source supplies no definition of pseudoflat and no proposed list of equivalent properties. In particular it is not replaced by flatness, a guessed Tor-vanishing condition, or a finitely-presented approximation property. The relevant side convention for pseudoflatness must also be fixed. This named-property gap prevents a self-contained characterization theorem for the second question.

The two requests are retained separately; the record does not assert that the cotorsion condition and the envelope-cokernel condition are themselves equivalent. The supplied literature triage checked 2026-08-21 reports no authoritative complete resolution of the extracted pair, without a specific example, characterization, or pseudoflatness definition.
\subsection{Short English statement}
Seek a noncoherent ring whose cotorsion left modules are all pure-injective, and characterize rings whose every pure-injective-envelope cokernel is pseudoflat. The coherence side and the definition of pseudoflat are essential missing data.
\subsection{Sources}
\begin{itemize}
\item Local UnsolvedMath registry, permanent ID 30000678, catalogue entry, \emph{Cotorsion Modules and Pure-Injective Envelopes}.
\item \emph{Model Theory and Groups} (2007), Oberwolfach Reports, pp. 83--138, DOI \url{https://doi.org/10.4171/owr/2007/02}.
\item Supplied source-verification claim and literature triage checked 2026-08-21 cite only the report; no separate ring characterization, example, or pseudoflatness definition is supplied.
\end{itemize}
\end{document}
